\begin{lemma}[Front member existence]
Let $F$ be a front on $X$, and let $x$ be an increasing enumeration whose
range is contained in $X$.  Then there is an $s\in F$ with
$s\segs\operatorname{range}(x)$.
\end{lemma}
\begin{proof}
By \noderef{IncSeq}, $x$ is an order embedding of $\mathbb N$ into
$\mathbb N$, and hence $x$ is injective.  Let
$R=\operatorname{range}(x)$.  For every $N\in\mathbb N$, the $N+1$
values
\[
x(0),x(1),\ldots,x(N)
\]
belong to $R$ and are pairwise distinct, because equality of two of them
would contradict injectivity.  Thus $R$ has at least $N+1$ elements for
every $N$.  If $R$ were finite, let $N$ be its cardinality.  The displayed
$N+1$ pairwise distinct elements of $R$ would then give more elements than
$R$ has, a contradiction.  Therefore $R$ is infinite.

The hypothesis $h_x$ says exactly that $R\subseteq X$.  Hence the density
clause in the definition of a front, \noderef{Front}, applied to the
infinite set $R\subseteq X$, yields an $s\in F$ such that
\[
\mathsf{ProperPrefixSet}(s,R).
\]
By the definition of \noderef{ProperPrefixSet}, this means that some
$n\in R$ satisfies
\[
s=\{k\in R\mid k<n\}.
\]
Since $R=\operatorname{range}(x)$, this says precisely that $s$ is the
finite initial segment of the range of $x$ ending immediately before the
member $n$ of that range.  Together with $s\in F$, the definition of
\noderef{FrontPrefix} therefore gives
\[
\mathsf{FrontPrefix}(F,s,x),
\]
which proves the required existential statement.
\end{proof}
